\textbf{1. Kiến thức cần nhớ}
\begin{itemize}
\item \textbf{Đơn thức} là biểu thức đại số chỉ gồm một số, hoặc một biến, hoặc một tích giữa các số và các biến.
\begin{itemize}
\item Số /bin/zsh$ được gọi là \textit{đơn thức không}.
\item Biểu thức chứa phép cộng, trừ giữa các biến hoặc chứa biến ở mẫu (như $\dfrac{1}{x}$, $\dfrac{x+y}{2}$) \textbf{không} phải là đơn thức.
\end{itemize}
\item \textbf{Đơn thức thu gọn} là đơn thức chỉ gồm tích của một số với các biến, mà mỗi biến đã được nâng lên lũy thừa với số mũ nguyên dương và chỉ xuất hiện một lần.
\begin{itemize}
\item Số thực khác /bin/zsh$ gọi là \textit{phần hệ số}.
\item Tích các biến gọi là \textit{phần biến}.
\item Khi viết đơn thức thu gọn, ta viết hệ số trước, phần biến sau và các biến theo thứ tự chữ cái.
\end{itemize}
\item \textbf{Bậc của đơn thức}:
\begin{itemize}
\item Bậc của đơn thức có hệ số khác /bin/zsh$ là \textbf{tổng số mũ} của tất cả các biến có trong đơn thức đó.
\item Số thực khác /bin/zsh$ là đơn thức bậc /bin/zsh$.
\item Số /bin/zsh$ là đơn thức không có bậc.
\end{itemize}
\item \textbf{Đơn thức đồng dạng} là hai đơn thức có hệ số khác /bin/zsh$ và có cùng phần biến.
\item \textbf{Cộng, trừ đơn thức đồng dạng}: Để cộng (hay trừ) các đơn thức đồng dạng, ta cộng (hay trừ) các hệ số với nhau và giữ nguyên phần biến:
80493\boxed{A x^m y^n \pm B x^m y^n = (A \pm B) x^m y^n}80493
\end{itemize}

\textbf{2. Phương pháp giải các dạng}

%% DANG 8D1?1-1
\textbf{Dạng 1. Nhận biết đơn thức, đơn thức thu gọn, bậc của đơn thức}
\begin{enumerate}
\item \textbf{Nhận biết đơn thức}: Kiểm tra biểu thức chỉ chứa phép nhân, lũy thừa của các số và biến (không chứa phép cộng, trừ giữa các biến, không chứa biến ở mẫu).
\item \textbf{Thu gọn đơn thức}:
\begin{itemize}
\item Nhân các hệ số số học với nhau.
\item Nhân các lũy thừa của cùng một biến: ^m \cdot x^n = x^{m+n}$.
\end{itemize}
\item \textbf{Tìm bậc}: Sau khi đã thu gọn, cộng số mũ của tất cả các biến lại.
\end{enumerate}
Lưu ý: Nếu một biến không ghi số mũ thì hiểu số mũ của nó là $ (ví dụ:  = x^1$).

%% DANG 8D1?1-2
\textbf{Dạng 2. Đơn thức đồng dạng; cộng, trừ đơn thức đồng dạng}
\begin{enumerate}
\item \textbf{Nhận biết đồng dạng}: Thu gọn các đơn thức về dạng chuẩn, kiểm tra xem hệ số có khác /bin/zsh$ và phần biến có giống nhau hoàn toàn hay không.
\item \textbf{Thực hiện phép tính}:
\begin{itemize}
\item Đặt phần biến chung ra ngoài (hoặc giữ nguyên phần biến).
\item Mở ngoặc tính tổng hoặc hiệu các hệ số.
\end{itemize}
\end{enumerate}
Lưu ý: Các đơn thức chỉ khác nhau về thứ tự các biến vẫn là đồng dạng (ví dụ: xy^2$ và y^2x$ là đồng dạng).

%% DANG 8D1?1-3
\textbf{Dạng 3. Tính giá trị của đơn thức}
\begin{enumerate}
\item \textbf{Bước 1}: Thu gọn đơn thức (nếu đơn thức chưa thu gọn).
\item \textbf{Bước 2}: Thay các giá trị cho trước của biến vào đơn thức đã thu gọn.
\item \textbf{Bước 3}: Thực hiện các phép tính theo thứ tự: lũy thừa $\rightarrow$ nhân, chia.
\end{enumerate}
Lưu ý: Khi thay số âm hoặc phân số, cần đặt trong dấu ngoặc đơn, ví dụ: ^3 = -8$, $\left(-\dfrac{1}{2}\right)^2 = \dfrac{1}{4}$.
